\documentclass[10pt,a4paper]{amsart}
\usepackage{amsmath,amssymb,amsthm,graphicx}
\usepackage{verbatim}
\newenvironment{reference}{\comment}{\endcomment}
%\newenvironment{reference}{}{}
\newenvironment{slogan}{\comment}{\endcomment}
\newenvironment{history}{\comment}{\endcomment}

% For commutative diagrams we use Xy-pic
\usepackage[all]{xy}

% We use 2cell for 2-commutative diagrams.
\xyoption{2cell}
% Theorem environments.
%
\theoremstyle{plain}
\newtheorem{theorem}[subsection]{Theorem}
\newtheorem{proposition}[subsection]{Proposition}
\newtheorem{lemma}[subsection]{Lemma}

\theoremstyle{definition}
\newtheorem{definition}[subsection]{Definition}
\newtheorem{example}[subsection]{Example}
\newtheorem{exercise}[subsection]{Exercise}
\newtheorem{situation}[subsection]{Situation}

\theoremstyle{remark}
\newtheorem{remark}[subsection]{Remark}
\newtheorem{remarks}[subsection]{Remarks}

\numberwithin{equation}{subsection}

% Macros
%
\def\lim{\mathop{\mathrm{lim}}\nolimits}
\def\colim{\mathop{\mathrm{colim}}\nolimits}
\def\Spec{\mathop{\mathrm{Spec}}}
\def\Hom{\mathop{\mathrm{Hom}}\nolimits}
\def\Ext{\mathop{\mathrm{Ext}}\nolimits}
\def\SheafHom{\mathop{\mathcal{H}\!\mathit{om}}\nolimits}
\def\SheafExt{\mathop{\mathcal{E}\!\mathit{xt}}\nolimits}
\def\Sch{\mathit{Sch}}
\def\Mor{\mathop{\mathrm{Mor}}\nolimits}
\def\Ob{\mathop{\mathrm{Ob}}\nolimits}
\def\Sh{\mathop{\mathit{Sh}}\nolimits}
\def\NL{\mathop{N\!L}\nolimits}
\def\CH{\mathop{\mathrm{CH}}\nolimits}
\def\proetale{{pro\text{-}\acute{e}tale}}
\def\etale{{\acute{e}tale}}
\def\QCoh{\mathit{QCoh}}
\def\Ker{\mathop{\mathrm{Ker}}}
\def\Im{\mathop{\mathrm{Im}}}
\def\Coker{\mathop{\mathrm{Coker}}}
\def\Coim{\mathop{\mathrm{Coim}}}

% Boxtimes
%
\DeclareMathSymbol{\boxtimes}{\mathbin}{AMSa}{"02}

%
% Macros for moduli stacks/spaces
%
\def\QCohstack{\mathcal{QC}\!\mathit{oh}}
\def\Cohstack{\mathcal{C}\!\mathit{oh}}
\def\Spacesstack{\mathcal{S}\!\mathit{paces}}
\def\Quotfunctor{\mathrm{Quot}}
\def\Hilbfunctor{\mathrm{Hilb}}
\def\Curvesstack{\mathcal{C}\!\mathit{urves}}
\def\Polarizedstack{\mathcal{P}\!\mathit{olarized}}
\def\Complexesstack{\mathcal{C}\!\mathit{omplexes}}
% \Pic is the operator that assigns to X its picard group, usage \Pic(X)
% \Picardstack_{X/B} denotes the Picard stack of X over B
% \Picardfunctor_{X/B} denotes the Picard functor of X over B
\def\Pic{\mathop{\mathrm{Pic}}\nolimits}
\def\Picardstack{\mathcal{P}\!\mathit{ic}}
\def\Picardfunctor{\mathrm{Pic}}
\def\Deformationcategory{\mathcal{D}\!\mathit{ef}}
\usepackage{cleveref}
\crefname{originalreference}{original reference}{original references}
\Crefname{originalreference}{Original reference}{Original references}
\makeatletter
\newlabel{lemma-faithful-general}{{0BLN}{source}{}{sourcepointer-1}{}}
\newlabel{lemma-faithful-general@cref}{{[originalreference][0][]0BLN}{[1][0][]source}{}{}{}}
\makeatother
\crefname{originallemma}{original lemma}{original lemmas}
\Crefname{originallemma}{Original lemma}{Original lemmas}
\makeatletter
\newlabel{spaces-pushouts-lemma-coequalizer}{{0AGG}{source}{}{sourcepointer-2}{}}
\newlabel{spaces-pushouts-lemma-coequalizer@cref}{{[originallemma][0][]0AGG}{[1][0][]source}{}{}{}}
\makeatother
\crefname{originallemma}{original lemma}{original lemmas}
\Crefname{originallemma}{Original lemma}{Original lemmas}
\makeatletter
\newlabel{spaces-pushouts-lemma-essentially-constant}{{0AGK}{source}{}{sourcepointer-3}{}}
\newlabel{spaces-pushouts-lemma-essentially-constant@cref}{{[originallemma][0][]0AGK}{[1][0][]source}{}{}{}}
\makeatother
\crefname{originalremark}{original remark}{original remarks}
\Crefname{originalremark}{Original remark}{Original remarks}
\makeatletter
\newlabel{spaces-pushouts-remark-essentially-constant}{{0AGL}{source}{}{sourcepointer-4}{}}
\newlabel{spaces-pushouts-remark-essentially-constant@cref}{{[originalremark][0][]0AGL}{[1][0][]source}{}{}{}}
\makeatother
\crefname{originalsection}{original section}{original sections}
\Crefname{originalsection}{Original section}{Original sections}
\makeatletter
\newlabel{spaces-pushouts-section-coequalizer-glue}{{0AGF}{source}{}{sourcepointer-5}{}}
\newlabel{spaces-pushouts-section-coequalizer-glue@cref}{{[originalsection][0][]0AGF}{[1][0][]source}{}{}{}}
\makeatother
\title{Gluing morphisms along a proper modification is controlled by the formal completion}
\author{The Stacks Project Authors}
\date{stacks-project:restricted.tex:local-snapshot-20261012; GNU FDL 1.2 or later}
\begin{document}
\maketitle
\noindent\textit{Editorial note.} All delivered mathematical text is the author's own, copied from the pinned Stacks Project chapter file restricted.tex; no statement or proof step was rewritten, reordered or removed.\par\medskip
\noindent\textit{Editorial note.} The counted proof is the author's complete argument: a uniqueness step via the gluing lemma, an etale-descent reduction to the affine Noetherian case, an equality over the self fibre product, and a coequalizer-system construction finished by essential constancy, all inside the delivered span\par\medskip
\noindent\textit{Editorial note.} The label proposition-glue-modification names the counted statement, and the proof environment immediately following it in the source is delivered verbatim as its proof.\par\medskip
\noindent\textit{Editorial note.} The counted statement is carried by the author's own proposition environment in the chapter source.\par\medskip
\noindent\textit{Allowed standard prerequisites (editorial).}\begin{itemize}
\item The formal completion X\_\{/T\} of an algebraic space along a closed subset, and its functoriality under the modification f : X' to X which is an isomorphism over the complementary open
\item The faithful gluing lemma of the same chapter (equality of morphisms from agreement on the open and after completion)
\item The coequalizer construction in the pushouts section of Pushouts of Spaces, including essential constancy of the system X\_n with value X
\end{itemize}\medskip
\ifdefined\hypertarget\hypertarget{sourcepointer-1}{}\fi
\noindent\textit{Source reference.} Original reference 0BLN: sources/stacks-restricted/pinned-original-source.tex, source lines 6484--6484. Prerequisite pointer invoked at proof line(s) 6657; it names a standard fact of the same chapter that the counted argument uses, and every step of the counted proof built on it is retained verbatim in the delivered body.\par\medskip
\ifdefined\hypertarget\hypertarget{sourcepointer-2}{}\fi
\noindent\textit{Source reference.} Original lemma 0AGG: sources/stacks-chow/tags.txt, source lines 13194--13194. Prerequisite pointer invoked at proof line(s) 6698; it names a standard fact of a sibling chapter that the counted argument uses, and every step of the counted proof built on it is retained verbatim in the delivered body.\par\medskip
\ifdefined\hypertarget\hypertarget{sourcepointer-3}{}\fi
\noindent\textit{Source reference.} Original lemma 0AGK: sources/stacks-chow/tags.txt, source lines 13198--13198. Prerequisite pointer invoked at proof line(s) 6713; it names a standard fact of a sibling chapter that the counted argument uses, and every step of the counted proof built on it is retained verbatim in the delivered body.\par\medskip
\ifdefined\hypertarget\hypertarget{sourcepointer-4}{}\fi
\noindent\textit{Source reference.} Original remark 0AGL: sources/stacks-chow/tags.txt, source lines 13199--13199. Prerequisite pointer invoked at proof line(s) 6714; it names a standard fact of a sibling chapter that the counted argument uses, and every step of the counted proof built on it is retained verbatim in the delivered body.\par\medskip
\ifdefined\hypertarget\hypertarget{sourcepointer-5}{}\fi
\noindent\textit{Source reference.} Original section 0AGF: sources/stacks-chow/tags.txt, source lines 13193--13193. Prerequisite pointer invoked at proof line(s) 6671; it names a standard fact of a sibling chapter that the counted argument uses, and every step of the counted proof built on it is retained verbatim in the delivered body.\par\medskip
\begin{proposition}
\label{proposition-glue-modification}
Let $S$ be a scheme. Let $X$ be a locally Noetherian algebraic space over $S$.
Let $T \subset |X|$ be a closed subset with complementary open subspace
$U \subset X$. Let $f : X' \to X$ be a proper morphism
of algebraic spaces such that $f^{-1}(U) \to U$ is an isomorphism.
For any algebraic space $W$ over $S$ the map
$$
\Mor_S(X, W) \longrightarrow
\Mor_S(X', W) \times_{\Mor_S(X'_{/T}, W)} \Mor_S(X_{/T}, W)
$$
is bijective.
\end{proposition}

\begin{proof}
Let $w' : X' \to W$ and $\hat w : X_{/T} \to W$ be morphisms
which determine the same morphism $X'_{/T} \to W$ by composition
with $X'_{/T} \to X$ and $X'_{/T} \to X_{/T}$. We have to prove
there exists a unique morphism $w : X \to W$ whose composition
with $X' \to X$ and $X_{/T} \to X$ recovers $w'$ and $\hat w$.
The uniqueness is immediate from Lemma \ref{lemma-faithful-general}.

\medskip\noindent
The assumptions on $T$ and $f$ are preserved by base change
by any \'etale morphism $X_1 \to X$ of algebraic spaces.
Since formal algebraic spaces are sheaves for the
\'etale topology and since we already have the uniqueness,
it suffices to prove existence after replacing $X$ by
the members of an \'etale covering. Thus we may assume $X$ is
an affine Noetherian scheme.

\medskip\noindent
Assume $X$ is an affine Noetherian scheme. We will construct
the morphism $w : X \to W$ using the material in
Pushouts of Spaces, Section \ref{spaces-pushouts-section-coequalizer-glue}.
It makes sense to read a little bit of the material
in that section before continuing the read the proof.

\medskip\noindent
Set $X'' = X' \times_X X'$ and consider the two morphisms
$a = w' \circ \text{pr}_1 : X'' \to W$ and 
$b = w' \circ \text{pr}_2 : X'' \to W$.
Then we see that $a$ and $b$ agree over the open $U$
and that $a_{/T}$ and $b_{a/T}$ agree (as these are
both equal to the composition $X''_{/T} \to X_{/T} \to W$
where the second arrow is $\hat w$).
Thus by Lemma \ref{lemma-faithful-general}
we see $a = b$.

\medskip\noindent
Denote $Z \subset X$ the reduced induced closed subscheme structure
on $T$. For $n \geq 1$ denote $Z_n \subset X$ the $n$th infinitesimal
neighbourhood of $Z$. Denote $w_n = \hat w|_{Z_n} : Z_n \to W$
so that we have $\hat w = \colim w_n$ on $X_{/T} = \colim Z_n$.
Set $Y_n = X' \amalg Z_n$. Consider the two projections
$$
s_n, t_n : R_n = Y_n \times_X Y_n \longrightarrow Y_n
$$
Let $Y_n \to X_n \to X$ be the coequalizer of $s_n$ and $t_n$ as in
Pushouts of Spaces, Section \ref{spaces-pushouts-section-coequalizer-glue}
(in particular this coequalizer exists, has good properties, etc, see
Pushouts of Spaces, Lemma \ref{spaces-pushouts-lemma-coequalizer}).
By the result $a = b$ of the previous parapgraph and the agreement
of $w'$ and $\hat w$ over $X'_{/T}$ we see that the morphism
$$
w' \amalg w_n : Y_n \longrightarrow W
$$
equalizes the morphisms $s_n$ and $t_n$. Hence we see that for all $n \geq 1$
there is a morphism $w^n : X_n \to W$ compatible with $w'$ and $w_n$.
Moreover, for $m \geq 1$ the composition
$$
X_n \to X_{n + m} \xrightarrow{w^{n + m}} W
$$
is equal to $w^n$ by construction (as the corresponding statement holds
for $w' \amalg w_{n + m}$ and $w' \amalg w_n$). By
Pushouts of Spaces, Lemma
\ref{spaces-pushouts-lemma-essentially-constant} and
Remark \ref{spaces-pushouts-remark-essentially-constant}
the system of algebraic spaces $X_n$ is essentially constant
with value $X$ and we conclude.
\end{proof}

\end{document}
